\documentclass[../thesis.tex]{subfiles}

\begin{document}
	\counterwithout{section}{chapter}
	\setcounter{section}{0}

	Описание разработанного в рамках проекта сервиса потоковой обработки данных.

	\section{Формулировка}

	Модуль предназначен для оценки эмиттанса и энергетического разброса пучка заряженных частиц. По
	своей сути модуль выполняет численную минимизацию заданной физической модели с использованием
	полученных измерений размера пучка в качестве наблюдений. Алгоритм разработан как для обработки данных
	в реальном времени, так и для воспроизведения эксперимента.

	\section{Физическое обоснование}

	\subsection{Описание физической задачи}

	Моделируется зависимость эмиттанса и энергетического разброса пучка заряженных частиц от его
	тока. Уравнения \ref{eq:1} и \ref{eq:2} описывают зависимость размеров пучка от параметров (см. \cite{binp}):

	\begin{align}
		\label{eq:1}\sigma_{x}^{2}= \beta_{x}\varepsilon + (D_{x}\sigma_{e})^{2}, \\
		\label{eq:2}\sigma_{z}^{2}= \beta_{z}\varepsilon + (D_{z}\sigma_{e})^{2},
	\end{align}

	\noindent
	где $\sigma_{x, z}$ -- размер пучка, $\beta_{x, z}$ -- бета-функция, $D_{x, z}$ -- дисперсия
	орбиты пучка, $\sigma_{e}$ -- энергетический разброс, $\varepsilon$ -- эмиттанс.

	\noindent
	Допущения:

	\begin{itemize}
		\item Размер пучка в накопительном кольце определяется поперечным распределением частиц,
			которое является гауссовским при низких токах (микроамперы). Таким образом, под размером
			пучка подразумевается ширина гауссовского распределения;

		\item Измерения размера пучка не синхронизированы по времени и требуют дополнительной предварительной
			обработки;

		\item Измерения размера пучка имеют определённую техническую точность, определяемую: фокусировкой
			изображения в плоскости CCD-сенсора, размером зерна CCD-сенсора, алгоритмом аппроксимации
			изображения пучка двумерной гауссовой функцией, переэкспонированием CCD-сенсора и
			наличием бликов.
	\end{itemize}

	\noindent
	С учётом последнего пункта, для практических целей приведённые выше выражения могут быть
	переписаны следующим образом:

	\begin{align}
		\sigma_{x}^{2}- \beta_{x}\varepsilon + (D_{x}\sigma_{e})^{2}\ge 0, \\
		\sigma_{z}^{2}- \beta_{z}\varepsilon + (D_{z}\sigma_{e})^{2}\ge 0.
	\end{align}
\end{document}